\subsection{Contador Síncrono MOD $10$}

Seja $D$ o estado atual do contador, restrito ao conjunto de valores válidos.

\paragraph{Incremento}
A operação de incremento é definida por:
\[
D \mapsto (D + 1) \bmod 10
\]

A Tabela~\ref{tab:inc_mod10} apresenta a função de transição.

\begin{table} [H] \label{tab:dec_mod10}
\begin{center}
\begin{tabular}{c|cc}
    $D$ & $O$ & Carry \\
    \hline
    $0000$ & $0001$ & $0$ \\
    $0001$ & $0010$ & $0$ \\
    $0010$ & $0011$ & $0$ \\
    $0011$ & $0100$ & $0$ \\
    $0100$ & $0101$ & $0$ \\
    $0101$ & $0110$ & $0$ \\
    $0110$ & $0111$ & $0$ \\
    $0111$ & $1000$ & $0$ \\
    $1000$ & $1001$ & $0$ \\
    $1001$ & $0000$ & $1$ \\
    $\vdots$ & $\vdots$ & $\vdots$ \\
\end{tabular}
\end{center}
\end{table}

As expressões booleanas minimizadas são:

\begin{align*}
    O_3 &= (\lnot D_3 \land D_2 \land D_1 \land D_0) \lor (D_3 \land \lnot D_2 \land \lnot D_1 \land \lnot D_0) \\[4pt]
    O_2 &= (\lnot D_3 \land \lnot D_2 \land D_1 \land D_0) \lor (\lnot D_3 \land D_2 \land \lnot D_1) \lor (\lnot D_3 \land D_2 \land \lnot D_0) \\[4pt]
    O_1 &= (\lnot D_3 \land \lnot D_1 \land D_0) \lor (\lnot D_3 \land D_1 \land \lnot D_0) \\[4pt]
    O_0 &= (\lnot D_3 \land \lnot D_0) \lor (\lnot D_2 \land \lnot D_1 \land \lnot D_0) \\[4pt]
    Carry &= D_3 \land \lnot D_2 \land \lnot D_1 \land D_0
\end{align*}

O sinal \textit{Carry} indica overflow.

As Figuras~\ref{fig:incmod10.png} e~\ref{fig:incmod10-red.png} apresentam, respectivamente, a implementação lógica do circuito e sua realização em redstone.

\imgbig{incmod10.png}{Circuito lógico para incremento em módulo 10.}
\imgbig{incmod10-red.png}{Implementação em redstone do circuito de incremento em módulo 10.}

\paragraph{Decremento}

A operação de decremento é definida por:
\[
D \mapsto (D - 1) \bmod 10
\]

A Tabela~\ref{tab:dec_mod10} apresenta a função de transição.

\begin{table} [H] \label{tab:dec_mod10}
\begin{center}
\begin{tabular}{c|cc}
    $D$ & $O$ & Carry \\
    \hline
    $0000$ & $1001$ & $1$ \\
    $0001$ & $0000$ & $0$ \\
    $0010$ & $0001$ & $0$ \\
    $0011$ & $0010$ & $0$ \\
    $0100$ & $0011$ & $0$ \\
    $0101$ & $0100$ & $0$ \\
    $0110$ & $0101$ & $0$ \\
    $0111$ & $0110$ & $0$ \\
    $1000$ & $0111$ & $0$ \\
    $1001$ & $1000$ & $0$ \\
    $\vdots$ & $\vdots$ & $\vdots$ \\
\end{tabular}
\end{center}
\end{table}

\begin{align*}
    O_3 &= (\lnot D_3 \land \lnot D_2 \land \lnot D_1 \land \lnot D_0) \lor (D_3 \land D_0) \\[4pt]
    O_2 &= (D_2 \land D_0) \lor (D_2 \land D_1) \lor (D_3 \land \lnot D_0) \\[4pt]
    O_1 &= (D_1 \land D_0) \lor (D_2 \land \lnot D_1 \land \lnot D_0) \lor (D_3 \land \lnot D_0) \\[4pt]
    O_0 &= \lnot D_0 \\[4pt]
    Carry &= \lnot D_3 \land \lnot D_2 \land \lnot D_1 \land \lnot D_0
\end{align*}

O sinal \textit{Carry} indica underflow.

As Figuras~\ref{fig:decmod10.png} e~\ref{fig:decmod10-red.png} mostram a implementação lógica e sua versão em redstone.

\imgbig{decmod10-red.png}{Implementação em redstone do circuito de decremento em módulo 10.}
\imgbig{decmod10.png}{Circuito lógico para decremento em módulo 10.}
